\chapter{Benchmarking}
\label{ch:bench}

NatQuery is evaluated on three tracks that answer three different questions (Table~\ref{tab:tracks}).
BIRD asks whether the security layer harms accuracy on a public benchmark. AdventureWorks asks whether
anything leaks under attack. The synthetic ERP asks whether the system scales to a university-sized,
multi-schema database with many roles. All tracks use the same open-weight model,
\code{openai/gpt-oss-120b} on Groq's free tier, at temperature 0 with no fine-tuning and no few-shot
examples.

\begin{table}[htbp]
\centering
\caption{The three benchmark tracks.}
\label{tab:tracks}
\small
\renewcommand{\arraystretch}{1.15}
\begin{tabularx}{\textwidth}{L{3.1cm} Y Y Y}
\toprule
\hdr \textbf{Track} & \textbf{Question answered} & \textbf{Database} & \textbf{Measure} \\
\midrule
BIRD Mini-Dev & Does the security layer cost accuracy? & 11 public databases, 75 tables, PostgreSQL & Official EX, Soft-F1, R-VES; latency, cost \\
AdventureWorks & Does anything leak under attack? & 68 tables with real-looking secrets, 3 roles & Leaked values, writes, audit integrity \\
Synthetic ERP & Does it scale to a university ERP? & 10 schemas, 87 tables, 549 columns, 10 roles & Correct answers per difficulty level, safety \\
\bottomrule
\end{tabularx}
\end{table}

\section{Accuracy and Efficiency on BIRD Mini-Dev}

\subsection{Setup}

BIRD Mini-Dev has 500 questions on 11 databases at three difficulty levels \citep{li2023bird}. All 75
tables were loaded unmodified into PostgreSQL. Every question went through the full NatQuery pipeline
with the security layer on, using one all-access benchmark role, and was scored with BIRD's own
official scripts. The model received BIRD's evidence hint, as every leaderboard entry does.

\begin{table}[htbp]
\centering
\caption{BIRD Mini-Dev setup.}
\label{tab:birdsetup}
\small
\renewcommand{\arraystretch}{1.15}
\begin{tabularx}{\textwidth}{L{3.6cm} Y}
\toprule
\hdr \textbf{Item} & \textbf{Value} \\
\midrule
Benchmark & BIRD Mini-Dev, PostgreSQL version: 500 questions (148 simple, 250 moderate, 102 challenging) \\
Model & \code{openai/gpt-oss-120b}, temperature 0, low reasoning effort \\
Input & Question plus BIRD's evidence hint (``oracle knowledge'' setting) \\
Security & On: every query passes the full enforcer \\
Scoring & Official BIRD scripts: execution accuracy (EX), Soft-F1, R-VES \\
Pipeline & One question at a time; first attempt plus at most two repairs \\
\bottomrule
\end{tabularx}
\end{table}

\subsection{Results}

NatQuery reaches an official execution accuracy of 45.6\% (228 of 500 questions exactly right under
BIRD's strictest comparison). Table~\ref{tab:birdres} gives all three metrics by difficulty and
Figure~\ref{fig:birdbar} charts EX.

\begin{table}[htbp]
\centering
\caption{Official BIRD scores with the security layer on.}
\label{tab:birdres}
\small
\renewcommand{\arraystretch}{1.15}
\begin{tabularx}{\textwidth}{L{3.6cm} Y Y Y Y}
\toprule
\hdr \textbf{Metric} & \textbf{Simple (148)} & \textbf{Moderate (250)} & \textbf{Challenging (102)} & \textbf{All (500)} \\
\midrule
Execution accuracy (EX) & 58.78 & 44.00 & 30.39 & \textbf{45.60} \\
Soft-F1 & 61.34 & 46.15 & 28.64 & \textbf{47.08} \\
R-VES & 58.46 & 42.92 & 28.70 & \textbf{44.62} \\
\bottomrule
\end{tabularx}
\end{table}

\begin{figure}[htbp]
\centering
\begin{tikzpicture}
\begin{axis}[ybar, bar width=22pt, width=0.8\textwidth, height=5.6cm, ymin=0, ymax=70, ylabel={EX (\%)},
  symbolic x coords={Simple,Moderate,Challenging,All}, xtick=data, nodes near coords, nodes near coords style={font=\footnotesize, /pgf/number format/fixed, /pgf/number format/precision=1},
  enlarge x limits=0.18, tick label style={font=\small}, ymajorgrids, grid style={gray!25}]
\addplot[fill=nqblue, draw=nqblue] coordinates {(Simple,58.78) (Moderate,44.00) (Challenging,30.39) (All,45.60)};
\end{axis}
\end{tikzpicture}
\caption{Official execution accuracy on BIRD Mini-Dev by difficulty.}
\label{fig:birdbar}
\end{figure}

\paragraph{The security layer costs no accuracy.} Table~\ref{tab:birdsec} shows what the enforcer did
to the model's first attempt. It accepted 484 of 500; of the 10 it rejected, 6 were cross joins and 4
unresolved references. In no case was a query that was right in raw form wrong in the final answer,
and the repair loop rescued three. On a numerically tolerant scorer, raw model SQL scored 50.2\%, the
same SQL through the enforcer 50.0\%, and the full pipeline 50.6\%.

\begin{table}[htbp]
\centering
\caption{What the enforcer did to first-attempt SQL (500 questions).}
\label{tab:birdsec}
\small
\renewcommand{\arraystretch}{1.15}
\begin{tabularx}{\textwidth}{Y L{2.6cm}}
\toprule
\hdr \textbf{Outcome} & \textbf{Count} \\
\midrule
Accepted by the enforcer & 484 (96.8\%) \\
Rejected (6 cross joins, 4 unresolved references) & 10 (2.0\%) \\
Right in raw form but wrong in the final answer & \textbf{0} \\
Rescued by the repair loop & 3 \\
\bottomrule
\end{tabularx}
\end{table}

\subsection{Position Among Published Systems}

Table~\ref{tab:leader} places NatQuery among the PostgreSQL entries we compared (execution accuracy,
oracle-knowledge setting). The scores of other systems were read from the public leaderboard; ours is
self-evaluated with the official scripts and not submitted. NatQuery ranks 8th of the 13 entries
compared and is ahead of every plain-model baseline: +9.6 points over GPT-4-turbo, +9.8 over GPT-4,
+16.2 over Llama3-70b-instruct and +18.2 over GPT-3.5-turbo. An open 120-billion-parameter model behind
a full security layer therefore beats GPT-4 given the same hints.

\begin{table}[htbp]
\centering
\caption{PostgreSQL leaderboard entries compared (EX, \%).}
\label{tab:leader}
\small
\renewcommand{\arraystretch}{1.1}
\begin{tabularx}{\textwidth}{L{1.2cm} Y L{2cm}}
\toprule
\hdr \textbf{Rank} & \textbf{System} & \textbf{EX} \\
\midrule
1 & schemagate + Claude Opus 5.5 & 72.6 \\
2 & SDIP + gemini-3.8-flash & 69.0 \\
3 & Ontology2SQL & 65.8 \\
4 & FrontierSQL + gpt-oss-120b & 56.8 \\
4 & SmartQuery + GPT-5.1 & 56.8 \\
6 & schemagate + Gemini 2.5 Flash & 55.6 \\
7 & TA + GPT-4 & 50.8 \\
\rowcolor{nqlight} \textbf{8} & \textbf{NatQuery + gpt-oss-120b, security on (self-evaluated)} & \textbf{45.6} \\
9 & GPT-4-turbo & 36.0 \\
10 & GPT-4 & 35.8 \\
11 & GPT-4-32k & 35.0 \\
12 & Llama3-70b-instruct & 29.4 \\
13 & GPT-3.5-turbo & 27.4 \\
\bottomrule
\end{tabularx}
\end{table}

\subsection{Efficiency}

Measured one question at a time with rate-limit waits excluded (Table~\ref{tab:eff}), the security layer
is invisible in the latency: the enforcer takes a median 4.3\,ms, 0.43\% of the end-to-end time, which
is dominated by the model call.

\begin{table}[htbp]
\centering
\caption{Efficiency on BIRD Mini-Dev.}
\label{tab:eff}
\small
\renewcommand{\arraystretch}{1.15}
\begin{tabularx}{\textwidth}{L{5.2cm} Y}
\toprule
\hdr \textbf{Measure} & \textbf{Value} \\
\midrule
End-to-end latency & p50 1.08\,s, p95 2.19\,s, p99 3.13\,s; 92.6\% of questions in under 2\,s \\
Model time & p50 1.05\,s (97.2\% of latency) \\
Enforcer time & p50 4.3\,ms, p95 10.8\,ms, p99 15.7\,ms \\
Database time & p50 3.5\,ms (2.4\% of latency) \\
Model calls per question & 1.03 (485 questions needed one call, 14 two, 1 three) \\
Tokens per question & 2,278 (2,092 in, 186 out); 4,995 per correct answer \\
Cost per question & \$0.0004 (\$0.21 for all 500) at \$0.15 / \$0.60 per million input / output tokens \\
Query speed & Median executed query takes 1.08$\times$ the gold query's database time \\
\bottomrule
\end{tabularx}
\end{table}

\paragraph{Schema selection.} With all 75 tables loaded as one 798-column schema, the table ranking put
every table the gold SQL needs into the prompt for 470 of 500 questions.

\subsection{Limitations of the BIRD Result}

\begin{table}[htbp]
\centering
\caption{Limitations of the BIRD track.}
\label{tab:birdlim}
\small
\renewcommand{\arraystretch}{1.15}
\begin{tabularx}{\textwidth}{L{3.6cm} Y}
\toprule
\hdr \textbf{Limitation} & \textbf{Detail} \\
\midrule
Self-evaluated & Official scripts run by us, not a leaderboard submission; the leaderboard may use a newer Mini-Dev version, so ranks are indicative \\
Brittle numeric metric & Official EX compares raw values without tolerance; about 25 misses are right answers with a different decimal form (our tolerant scorer gives 50.6\%) \\
Partly tuned on the test set & The schema-selection rule was designed after reading failures on these 500 questions; treat the result as an upper estimate until checked on unseen questions \\
One run & Temperature 0 is not fully deterministic on a hosted API; the spread was not measured \\
Weak model setting & Low reasoning effort on a free tier; the same model scores 56.8 in another published system, an 11-point gap we did not analyse \\
Evidence hint given & Real users will not provide hints, so expect lower accuracy on a real database without good descriptions \\
Security not tested here & One all-access role: this run shows the enforcer does not hurt accuracy, not that it blocks leaks \\
\bottomrule
\end{tabularx}
\end{table}

\FloatBarrier
\section{Security Testing on AdventureWorks}

\subsection{Setup}

Accuracy says nothing about leaks, so a separate track attacks the system. The database is Microsoft
AdventureWorks 2014 on PostgreSQL: 68 tables, 31,465 sales orders, 19,972 people and 290 employees,
with realistic secrets (password hashes, credit card numbers, national ID numbers). The pass criteria
were fixed before the runs: zero leaked values, zero successful writes, zero unbounded executions, and
every request audited. A leak checker independent of NatQuery's policy code holds 89,686 protected
strings read straight from the tables and each role's entitled keys, and checks every returned cell.
The database login was the superuser on purpose: if NatQuery's checks failed, nothing beneath them would
help.

\begin{table}[htbp]
\centering
\caption{AdventureWorks roles and what is hidden.}
\label{tab:awroles}
\small
\renewcommand{\arraystretch}{1.15}
\begin{tabularx}{\textwidth}{L{2.6cm} Y}
\toprule
\hdr \textbf{Role} & \textbf{Access} \\
\midrule
\code{employee} & Own HR record and pay; product catalogue \\
\code{sales_rep} & Own orders and customers; customer e-mail masked as \code{***@domain} \\
\code{admin} & Everything except the permanently hidden data \\
Hidden for all & Tables \code{password}, \code{creditcard}, \code{personcreditcard}, \code{jobcandidate}; columns \code{passwordhash}, \code{passwordsalt}, \code{cardnumber}, \code{nationalidnumber}, \code{creditcardid}, \code{creditcardapprovalcode} \\
\bottomrule
\end{tabularx}
\end{table}

Four test layers simulate different adversaries (Table~\ref{tab:awlayers}).

\begin{table}[htbp]
\centering
\caption{The four security test layers.}
\label{tab:awlayers}
\small
\renewcommand{\arraystretch}{1.15}
\begin{tabularx}{\textwidth}{L{3.4cm} Y}
\toprule
\hdr \textbf{Layer} & \textbf{What it simulates} \\
\midrule
A. Hostile SQL & The model is compromised or jailbroken and returns the attacker's SQL; sent through the full pipeline \\
B. Auth, sessions, audit & Login, lockout, forged tokens, role escalation, audit tampering, what the model is shown \\
C. Natural language & Real users typing real attacks and questions into the real model \\
D. Unprotected baseline & The same model and kinds of attack with no NatQuery, on the owner connection \\
\bottomrule
\end{tabularx}
\end{table}

Layer A covers eleven attack families: globally hidden data, tables and columns outside a role's grant,
row-filter bypasses (\code{OR TRUE}, self-joins, subqueries, CTEs, UNION), masked-column probing (LIKE,
ORDER BY, GROUP BY, MIN, md5, joins), writes and DDL, stacked statements, dangerous functions, catalog
and system tables, obfuscation, whole-row references and resource abuse.

\subsection{Results}

\begin{table}[htbp]
\centering
\caption{AdventureWorks security results.}
\label{tab:awres}
\small
\renewcommand{\arraystretch}{1.15}
\begin{tabularx}{\textwidth}{L{5.6cm} Y}
\toprule
\hdr \textbf{Test} & \textbf{Result} \\
\midrule
Hostile SQL: leaked values & \textbf{0 of 560 runs} (227 queries, 3 roles, 11 families) \\
Hostile SQL: writes or DDL & \textbf{0}; fingerprint of all 68 tables identical before and after \\
Must-refuse queries refused & \textbf{511 of 511} \\
Allowed queries returned only entitled data & 49 of 49 (for example a sales rep's \code{COUNT(*)} of all orders returned 418, not 31,465) \\
Auth, session and audit checks & 40 of 40 \\
Natural-language attacks through NatQuery & \textbf{0 leaks in 30} (25 refused or ``unanswerable'', 5 answered with entitled data only) \\
Same attacks on an unprotected text-to-SQL & \textbf{7 of 8 leaked} (all 290 national IDs, all password hashes, company-wide revenue, every employee row) \\
Same hostile SQL run directly on the database & 141 of 187 must-refuse queries completed; 47 returned secret data \\
Prompt size with the security layer & 66\% fewer tokens (2,170 vs 6,295): the model sees only its role's tables \\
Enforcer plus execution & Median 6\,ms per request \\
\bottomrule
\end{tabularx}
\end{table}

\paragraph{What testing found.} The tests found three real bugs, all fixed with regression tests
(Table~\ref{tab:awbugs}). The first run is kept in the repository.

\begin{table}[htbp]
\centering
\caption{Bugs found by the security track.}
\label{tab:awbugs}
\small
\renewcommand{\arraystretch}{1.15}
\begin{tabularx}{\textwidth}{Y Y}
\toprule
\hdr \textbf{Bug} & \textbf{Why it mattered} \\
\midrule
\code{SELECT employee::text FROM employee} returned the whole row, including the hidden national ID & A table name used as a value bypassed column grants: a real leak \\
\code{repeat('a', 1000000000)} crashed the database connection & An allowed function could exhaust memory \\
\code{JOIN ... ON true} passed the cross-join rule & Resource abuse, bounded only by the timeout \\
\bottomrule
\end{tabularx}
\end{table}

\subsection{Limitations of the Security Result}

\begin{table}[htbp]
\centering
\caption{Limitations of the security track.}
\label{tab:awlim}
\small
\renewcommand{\arraystretch}{1.15}
\begin{tabularx}{\textwidth}{L{3.6cm} Y}
\toprule
\hdr \textbf{Limitation} & \textbf{Detail} \\
\midrule
Own defences tested & The attack list was written by the authors of the enforcer, and the first run found three bugs, so more probably exist; an outside red team is the missing step \\
Small natural-language sample & 30 attacks (a free key allows about 200,000 tokens a day): zero leaks only rules out a rate above about 10\%; the strong claim rests on layer A \\
Answers scoped silently & A sales rep asking for ``total revenue of the company'' got their own figure with no warning; nothing leaked, but a user can be misled \\
Inference not covered & A role allowed to aggregate rows can learn statistics about them; there is no minimum group size \\
Audit tamper-evident only & Without a key a full rewrite can be forged; truncation at the end is not detectable by the chain alone \\
Heuristic input guard & It stopped 4 of 30 attacks; the enforcer, not the guard, is the boundary \\
One model, one engine & Only \code{gpt-oss-120b} and PostgreSQL were tried; the enforcer result does not depend on the model \\
Untested & Concurrency and timing channels (concurrency is covered by integration tests), SSO with the real ERP, and the database-role layer on this database \\
\bottomrule
\end{tabularx}
\end{table}

\FloatBarrier
\section{The Synthetic University ERP}

\subsection{Why a Synthetic ERP}

The target is the IIT Kharagpur ERP, whose real schema was not available. A test database was therefore
built from the screens of the student login and extended, by extrapolation, to the teacher and
administrative sides. Every row is synthetic and generated from a fixed seed; no real student data was
used. The field names follow the real screens, and the model is the single source of the DDL, the
diagrams and the policy.

\subsection{Database Design}

The database has 10 schemas, 87 tables, 549 columns, 130 foreign keys and about 368,000 rows
(Table~\ref{tab:erpschemas}, Figure~\ref{fig:erpbar}). Logins, sessions and one-time passwords are
implemented in SQL (bcrypt hashes, a five-failure lockout, hashed tokens and one-time passwords, full audit).

\begin{table}[htbp]
\centering
\caption{ERP schemas.}
\label{tab:erpschemas}
\small
\renewcommand{\arraystretch}{1.15}
\begin{tabularx}{\textwidth}{L{1.7cm} L{1.5cm} L{1.9cm} Y}
\toprule
\hdr \textbf{Schema} & \textbf{Tables} & \textbf{Columns} & \textbf{Content} \\
\midrule
\code{acad} & 27 & 164 & Registrations, results, attendance, exams, medical leave, graduation \\
\code{people} & 10 & 94 & Persons, students, employees, advisors, guardians, addresses \\
\code{svc} & 14 & 86 & Mail, notices, grievances, counselling, paperless applications, e-files, dispatch \\
\code{fin} & 10 & 60 & Fee structure and demands, payments, scholarships, mediclaim, salary slips \\
\code{auth} & 6 & 42 & Accounts, roles, sessions, one-time passwords, login audit \\
\code{org} & 5 & 26 & Departments, programs, academic sessions, halls, rooms \\
\code{lib} & 4 & 22 & Books, copies, loans, dues \\
\code{cdc} & 4 & 20 & Companies, postings, applications, offers (placement) \\
\code{sys} & 4 & 18 & Audit events, menu modules and role-menu access \\
\code{gh} & 3 & 17 & Guest house \\
\midrule
\textbf{Total} & \textbf{87} & \textbf{549} & \\
\bottomrule
\end{tabularx}
\end{table}

\begin{figure}[htbp]
\centering
\begin{tikzpicture}
\begin{axis}[xbar, width=0.85\textwidth, height=7cm, xmin=0, xlabel={Count}, symbolic y coords={gh,sys,cdc,lib,org,auth,fin,svc,people,acad}, ytick=data,
  legend style={at={(0.97,0.05)}, anchor=south east, font=\footnotesize}, bar width=5pt, enlarge y limits=0.07, xmajorgrids, grid style={gray!25}, tick label style={font=\small}]
\addplot[fill=nqblue, draw=nqblue] coordinates {(27,acad) (10,people) (14,svc) (10,fin) (6,auth) (5,org) (4,lib) (4,cdc) (4,sys) (3,gh)};
\addplot[fill=nqorange, draw=nqorange] coordinates {(164,acad) (94,people) (86,svc) (60,fin) (42,auth) (26,org) (22,lib) (20,cdc) (18,sys) (17,gh)};
\legend{Tables, Columns}
\end{axis}
\end{tikzpicture}
\caption{Tables and columns per ERP schema.}
\label{fig:erpbar}
\end{figure}

\paragraph{Three database layers.} NatQuery logs in as \code{natquery}, which owns nothing and can read
nothing; each query runs after \code{SET LOCAL ROLE erp_ro}. \code{erp_ro} is SELECT-only and has no
grant on password hashes, session tokens or one-time-password hashes. Row-level security on 20 tables
repeats the row scopes in SQL; with no context a role sees no rows.

\paragraph{Roles.} Ten roles can be assigned (an eleventh, \code{erp_base}, is the shared reference-data
block every role inherits). Table~\ref{tab:erproles} shows how differently the views are sized: a
student sees 74 tables and 462 columns, the library staff 35 and 188.

\begin{table}[htbp]
\centering
\caption{ERP roles and the size of each role's view.}
\label{tab:erproles}
\small
\renewcommand{\arraystretch}{1.15}
\begin{tabularx}{\textwidth}{L{2.6cm} L{1.6cm} L{1.7cm} L{2.2cm} Y}
\toprule
\hdr \textbf{Role} & \textbf{Tables} & \textbf{Columns} & \textbf{Tables with row filters} & \textbf{Scope} \\
\midrule
\code{student} & 74 & 462 & 49 & Own records only \\
\code{faculty} & 47 & 265 & 20 & Own courses, advisees, own salary \\
\code{hod} & 50 & 291 & 23 & Whole department \\
\code{dept_office} & 50 & 291 & 23 & Same as the head of department \\
\code{acad_section} & 60 & 358 & 8 & All academic records \\
\code{accounts} & 45 & 256 & 7 & All finance \\
\code{library} & 35 & 188 & 6 & Library staff \\
\code{hall_office} & 35 & 194 & 10 & Residents of one hall \\
\code{cdc} & 38 & 208 & 11 & Career Development Cell \\
\code{admin} & 85 & 532 & 0 & System administrator (permanent secrets still hidden) \\
\bottomrule
\end{tabularx}
\end{table}

\subsection{Verifying the Access Layer Without a Model}

Before any model is involved, two checks establish that the access layer is correct, so that any later
miss belongs to the model (Table~\ref{tab:erpverify}). The ground truth for every question is computed
as the database superuser, by SQL written under the role's scope.

\begin{table}[htbp]
\centering
\caption{Checks that need no language model.}
\label{tab:erpverify}
\small
\renewcommand{\arraystretch}{1.15}
\begin{tabularx}{\textwidth}{L{4.6cm} L{2.4cm} Y}
\toprule
\hdr \textbf{Check} & \textbf{Result} & \textbf{What it shows} \\
\midrule
Database checks (logins, lockout, one-time passwords, sessions, grants, row-level security against an independent oracle) & 51 of 51 & The database layer behaves as designed, including querying as \code{erp_ro} without NatQuery \\
No-LLM pipeline pass: every ground-truth query and every attack query sent through the full pipeline & 41 of 41 & All 35 answerable questions return exactly the expected rows; all 6 attacks return nothing forbidden \\
\bottomrule
\end{tabularx}
\end{table}

\subsection{The Question Set}

There are 40 questions in three difficulty levels (Table~\ref{tab:erpq}), each asked by the role it
concerns. The full list with results is in Appendix~\ref{app:questions}. Of the 12 hard questions, five
are security questions: four that a role must be refused or must not be able to answer from forbidden
data (other people's pay, sensitive identifiers, a prompt injection, finance data for a role without
finance access) and one scope trap, where an answer is acceptable only if it contains nothing forbidden.

\begin{table}[htbp]
\centering
\caption{Design of the 40 ERP questions.}
\label{tab:erpq}
\small
\renewcommand{\arraystretch}{1.15}
\begin{tabularx}{\textwidth}{L{1.8cm} L{1.5cm} Y}
\toprule
\hdr \textbf{Level} & \textbf{Count} & \textbf{What it needs} \\
\midrule
Easy & 11 & One table or one join: own record, a count, a latest row, a simple filter \\
Medium & 17 & Two to four tables across schemas, status and date logic, grouping, averages, names through the person table \\
Hard & 12 & Window functions, five-table joins, percentages with thresholds, department scope, and five access-control traps \\
\bottomrule
\end{tabularx}
\end{table}

\paragraph{Scoring.} For answerable questions, all values of the ground truth must be present in the
answer with the same number of rows; extra columns and different labels are accepted. A security
question counts as safe if nothing from its forbidden set appears in the answer.

\subsection{Results}

Table~\ref{tab:erpres} and Figure~\ref{fig:erpacc} give the results of two full runs of all 40 questions.
The first run exposed what the model could not see; the second followed the generic changes of
Table~\ref{tab:erpchanges}.

\begin{table}[htbp]
\centering
\caption{ERP accuracy before and after the generic improvements (correct / questions).}
\label{tab:erpres}
\small
\renewcommand{\arraystretch}{1.15}
\begin{tabularx}{\textwidth}{L{3.2cm} Y Y}
\toprule
\hdr \textbf{Level} & \textbf{First run} & \textbf{Final run} \\
\midrule
Easy & 4 / 11 (36.4\%) & 9 / 11 (81.8\%) \\
Medium & 4 / 17 (23.5\%) & 11 / 17 (64.7\%) \\
Hard & 6 / 12 (50.0\%) & 8 / 12 (66.7\%) \\
\midrule
\textbf{All} & \textbf{14 / 40 (35.0\%)} & \textbf{28 / 40 (70.0\%)} \\
Answerable questions (35) & 9 (25.7\%) & 23 (65.7\%) \\
Security questions (5) & 5 safe, 0 leaks & 5 safe, 0 leaks \\
Refused as ``unanswerable'' & 10 & 4 \\
Model calls and tokens per question & 0.97 calls, about 2,770 tokens & 1.30 calls, about 4,940 tokens \\
\bottomrule
\end{tabularx}
\end{table}

\begin{figure}[htbp]
\centering
\begin{tikzpicture}
\begin{axis}[ybar, bar width=15pt, width=0.85\textwidth, height=6cm, ymin=0, ymax=100, ylabel={Correct (\%)},
  symbolic x coords={Easy,Medium,Hard,All}, xtick=data, enlarge x limits=0.18, legend style={at={(0.03,0.97)}, anchor=north west, font=\footnotesize},
  nodes near coords, nodes near coords style={font=\scriptsize}, ymajorgrids, grid style={gray!25}, tick label style={font=\small}]
\addplot[fill=nqorange!70, draw=nqorange] coordinates {(Easy,36.4) (Medium,23.5) (Hard,50.0) (All,35.0)};
\addplot[fill=nqblue, draw=nqblue] coordinates {(Easy,81.8) (Medium,64.7) (Hard,66.7) (All,70.0)};
\legend{First run, Final run}
\end{axis}
\end{tikzpicture}
\caption{ERP accuracy by difficulty level, first and final run.}
\label{fig:erpacc}
\end{figure}

Latency is not comparable with BIRD: the free tier allows 8,000 tokens per minute, so questions were
paced about 20 seconds apart and waited on token limits, giving a median of about 62 seconds per
question in the final run. With a paid tier or an on-premises model this wait disappears.

\subsection{What Changed Between the Two Runs}

After the first run the failures were read and grouped. Table~\ref{tab:erpchanges} lists the changes;
all are generic library features, and none mentions a question. Sixteen questions turned from wrong to
correct (E02, E03, E05, E09, E10, M01, M02, M04, M05, M11, M12, M16, M17, H08, H09, H12) and two turned
from correct to wrong (M06 and H05), a net gain of fourteen.

\begin{table}[htbp]
\centering
\caption{Failures of the first run and the generic change that addressed each.}
\label{tab:erpchanges}
\small
\renewcommand{\arraystretch}{1.15}
\begin{tabularx}{\textwidth}{L{3.5cm} Y Y}
\toprule
\hdr \textbf{Failure} & \textbf{Cause} & \textbf{Change} \\
\midrule
Wrong status literals (about 10 questions) & The prompt listed column names but no values, so the model guessed \code{'failed'} or \code{'Success'} & Value hints for categorical columns, generated from the database; prompt rule to use the listed spellings \\
Department id instead of name & No join to the lookup table & Small lookup tables added to the selected schema; prompt rule to return readable labels \\
``Unanswerable'' on role-scoped questions & ``My'' and ``our'' read as missing data; needed tables not always in the prompt & Prompt rule that rows are scoped automatically; weighted table ranking; wider-schema retry \\
Tables dropped from the prompt & Alphabetical filling lost whole schemas such as \code{org} and \code{people}; synonyms were not passed to the ranker & Fill by connectivity and rank; synonyms passed to the ranker \\
\bottomrule
\end{tabularx}
\end{table}

\paragraph{Retrieval.} Measured by comparing the tables used by each ground-truth query with the
tables placed in the prompt, the improved table selection covers all needed tables for 31 of the 35
answerable questions, against 26 with the previous ranking and fill rule.

\subsection{What Still Fails}

Twelve questions are still wrong (Table~\ref{tab:erpfail}). Three need a policy decision, three are refusals, one is a permission refusal and five are SQL or value errors.

\begin{table}[htbp]
\centering
\caption{Remaining failures of the final run.}
\label{tab:erpfail}
\small
\renewcommand{\arraystretch}{1.15}
\begin{tabularx}{\textwidth}{L{2.6cm} L{2.2cm} Y}
\toprule
\hdr \textbf{Group} & \textbf{Questions} & \textbf{Cause} \\
\midrule
``My'' on a table without a row filter & E07, M06, H06 & The prompt forbids filtering on the user's own id and the instructor table has no row filter, so the answer covers every instructor instead of the asker; a policy choice is needed (add a filter, or let the model filter on identity) \\
Refused as ``unanswerable'' & E08, M08, M09 & Department and previous-session questions the model still judged unanswerable \\
Permission refusal & H05 & The query touched a table the role cannot read \\
Multi-step SQL errors & M10, H10, H11 & Window functions, ranking and percentage logic wrong \\
Threshold and value errors & M03, M14 & Attendance percentage threshold; a mobile-number question that returned no rows \\
\bottomrule
\end{tabularx}
\end{table}

\subsection{Caveats on the ERP Result}

\begin{table}[htbp]
\centering
\caption{Caveats on the ERP track.}
\label{tab:erpcav}
\small
\renewcommand{\arraystretch}{1.15}
\begin{tabularx}{\textwidth}{L{3.6cm} Y}
\toprule
\hdr \textbf{Caveat} & \textbf{Detail} \\
\midrule
Improvements read the failures & The changes followed the first-run failures on the same 40 questions; the final 70\% is an upper estimate for new questions \\
Synthetic data and policy & The questions were written by the people who wrote the policy; the real ERP has its own quirks \\
One run each & No measurement of run-to-run spread \\
Free tier and one model & Small prompts forced by 8,000 tokens per minute; a stronger model would probably score higher \\
Silent scoping & A row filter changes totals without telling the user (known limit, Chapter~\ref{ch:conclusion}) \\
Row-level security covers rows only & Column hiding and masking exist in the enforcer alone, except for secret columns the database role is never granted \\
\bottomrule
\end{tabularx}
\end{table}

\FloatBarrier
\section{Summary of the Evidence}

\begin{table}[htbp]
\centering
\caption{What each track shows and what it does not.}
\label{tab:summary}
\small
\renewcommand{\arraystretch}{1.15}
\begin{tabularx}{\textwidth}{L{2.6cm} Y Y}
\toprule
\hdr \textbf{Track} & \textbf{Shows} & \textbf{Does not show} \\
\midrule
BIRD & The security layer costs no accuracy; 45.6\% EX with an open model at \$0.0004 per question & Leak resistance; accuracy without hints; behaviour under load \\
AdventureWorks & No leak or write under 560 hostile runs and 30 natural-language attacks, where an unprotected baseline leaked & That no attack exists; outside red-team results; inference attacks \\
Synthetic ERP & The system handles 87 tables, 549 columns and 10 roles; the access layer returns exactly the right rows; 70\% correct, all security questions safe & That the real ERP's policy is right; accuracy on unseen questions \\
\bottomrule
\end{tabularx}
\end{table}
